\section{Planteamiento del problema y justificación}
~\label{sec:problema}

\subsection{Situación problema}

La gestión de infraestructuras de red en entornos locales 
presenta una dependencia estructural respecto a consolas de 
administración estáticas que condiciona tanto la continuidad como la 
calidad de la supervisión. Durante los periodos en los que el 
administrador se desplaza, la capacidad de observar el 
comportamiento de los activos se interrumpe, y los eventos que 
ocurren en esos intervalos solo se conocen cuando sus efectos ya se 
manifestaron en forma de caídas de servicio o lentitud reportada por 
los usuarios. La información disponible para el diagnóstico, además, 
suele limitarse a registros crudos en archivos PCAP o logs en texto 
plano, sin un procesamiento que la transforme en métricas 
contextualizadas ni un inventario que permita distinguir los 
dispositivos legítimos de aquellos que no deberían estar conectados. 
El resultado es una administración que opera sobre información 
parcial y reacciona tarde.

\subsection{Árbol de causas y efectos}

El árbol del problema (Fig.~\ref{fig:arbol-problema}) identifica 
cuatro causas estructurales y cuatro efectos directos que orientan el 
diseño de la solución.

\noindent\textbf{C1. Dependencia de modelos centralizados y estáticos.} 
Las herramientas convencionales operan desde un punto fijo dentro de 
la red y presentan la información en formatos poco adecuados para 
dispositivos móviles o pantallas reducidas, lo que hace impráctica la 
supervisión remota.

\noindent\textbf{C2. Ausencia de captura distribuida por endpoint.} 
La captura desde un único punto de observación ofrece una visión 
agregada del tráfico, pero no permite correlacionar flujos anómalos 
con el dispositivo que los origina.

\noindent\textbf{C3. Falta de contextualización de las métricas.} 
El volumen capturado por las herramientas tradicionales no se 
transforma en indicadores accionables como salud del enlace, perfiles 
de consumo, tasas de retransmisión TCP o tiempos de respuesta DNS, de 
modo que la información existe pero permanece inaccesible en un 
formato que dificulta el diagnóstico.

\noindent\textbf{C4. Carencia de inventario dinámico clasificado.} 
Los escaneos activos detectan la presencia de dispositivos mediante 
direcciones IP y MAC, pero no clasifican el tipo de equipo ni su rol 
dentro de la red, con lo cual el administrador vigila una red cuya 
composición exacta desconoce.

Estas causas se traducen en cuatro efectos operativos: tiempos de 
detección prolongados ante anomalías y amenazas (E1); interrupciones 
en la supervisión durante los periodos de movilidad del administrador 
(E2); diagnósticos imprecisos basados en información incompleta (E3); 
y presencia no detectada de dispositivos no autorizados, con la 
consecuente exposición a riesgos internos sobre la confidencialidad, 
integridad y disponibilidad de la información (E4).

\begin{figure}[ht]
\centering
\includegraphics[width=\columnwidth]{arbol_problema.png}
\caption{Árbol de causas y efectos del problema.}
~\label{fig:arbol-problema}
\end{figure}

\subsection{Pregunta problematizadora}

\textit{¿Cómo abordar la visibilidad y control sobre el 
comportamiento del tráfico de red y el estado de los activos 
conectados en una infraestructura, considerando la necesidad de 
detectar anomalías, métricas y dispositivos de forma independiente 
de la ubicación física del administrador?}

\subsection{Justificación}

Las limitaciones descritas trascienden el plano de la comodidad 
operativa y afectan la continuidad de los servicios, la capacidad 
de respuesta del equipo responsable y la exposición a riesgos de 
seguridad interna. Cada intervalo sin supervisión amplía la ventana 
durante la cual una anomalía puede escalar hacia una afectación de 
mayor alcance, y cada diagnóstico realizado sobre información 
incompleta incrementa la probabilidad de ciclos de prueba y error 
que consumen tiempo y pueden agravar la situación original. En 
organizaciones donde la red sostiene la totalidad de los procesos 
de comunicación interna y acceso a la información, esta brecha 
representa un costo directo en términos de productividad, calidad 
del servicio y posturas de seguridad.

El abordaje del problema requiere articular conocimientos de 
distintas áreas de la ingeniería de sistemas: captura y análisis de 
tráfico desde la perspectiva de redes de datos, procesamiento y 
persistencia de volúmenes significativos de información desde la 
perspectiva de bases de datos y sistemas distribuidos, y exposición 
de resultados al administrador desde la perspectiva del desarrollo 
de software y la movilidad. La naturaleza integradora del proyecto 
justifica su desarrollo como una propuesta unificada más que como 
la mejora aislada de un componente, y habilita un escenario de 
aplicación para las competencias de infraestructura tecnológica, 
seguridad de la información e ingeniería de software establecidas 
en el marco académico del Proyecto Integrador.